\section{Operatory Hilberta-Schmidta}

Rozważamy unitarne reprezentacje grup zwartych $G$ z unormowaną miarą Haara $dx$, to znaczy $dx(G)=1$.

% \begin{definition}{H-*-algebra}{}
%   Algebra $A$ z inwolucją $x\mapsto x^*$ która jest przestrzenią Hilberta z iloczynem skalarnym $\langle \cdot,\cdot\rangle$ nazywa się \buff{H-*-algebrą} jeśli 
%   \begin{itemize}
%     \item $(\forall\;x,y,z\in A)\;\langle xy, z\rangle=\langle y, x^*z\rangle$
%     \item $\|x^*\|=\|x\|$
%   \end{itemize}
%   Dodatkowo chcemy, żeby ta algebra była \acc{półprosta}, czyli dla $x\neq 0$ wymagamy $xx^*\neq 0$.
%     %
%     % Powiemy, że przestrzeń Hilberta $A$ z unitarnym działaniem $*$ nazywamy *-algebrą jeśli
%     % \begin{itemize}
%     %   \item $(\forall\;x,y,z\in A)\;\langle xy, z\rangle=\langle y, x^*z\rangle$
%     %   \item $\|x^*\|=\|x\|$
%     % \end{itemize}
% \end{definition}
%
% \begin{definition}{algebra Hilberta}{}
%   W literaturze algebrą Hilberta nazywamy algebrę $A$ z inwolucją $x\mapsto x^*$ oraz iloczynem skalarnym $\langle \cdot,\cdot\rangle$ taką, że
%   \begin{itemize}
%     \item $(\forall\;x,y,z\in A)\;\langle xy, z\rangle=\langle y, x^*z\rangle$
%     \item $(\forall\;x,y\in A)\;\langle x,y\rangle=\langle y^*,x^*\rangle$
%     \item dla dowolnego ustalonego $a\in A$ operator $x\mapsto ax$ jest ograniczony
%     \item przestrzeń liniowa rozpięta przez $xy$ dla $x,y\in A$ jest gęstym podzbiorem $A$.
%   \end{itemize}
% \end{definition}

Niech $T:\Hh\to\Hh$ będzie operatorem ograniczonym na przestrzeni Hilberta $\Hh$ z bazą $\{e_i\;:\;i\in I\}$. Definiujemy normę Hilberta-Schmidta operatora $T$ jako
$$\|T\|_{HS}^2:=\sum_{i\in I}\|Te_i\|_H^2.$$
Powiemy, że $T$ jest \buff{operatorem Hilberta-Schmidta} jeśli ma skończoną normę Hilberta-Schmidta.

\begin{fact}{}{}
  Norma Hilbera-Schmidta jest dobrze określona, tzn. nie zależy od wyboru bazy przestrzeni $\Hh$.
\end{fact}

\begin{proof}
  Zaczynamy od pokazania, że $\|T\|_{HS}^2=\|T^*\|_{HS}^2$. Pitagoras mówi nam, że jeśli $\{e_i\;:\;i\in I\}$ jest bazą ortonormalną $\Hh$, to
  $$\|x\|^2=\sum_i|\langle x, e_i\rangle|^2=\sum_i|\langle e_i, x\rangle|^2,$$
  bo $x=\sum\langle x,e_i\rangle e_i$. Stąd
  $$\|T\|_{HS}^2=\sum_i\|Te_i\|^2_{HS}=\sum_i\sum_j|\langle Te_i, e_j\rangle|^2=\sum_i\sum_j|\langle e_i, T^*e_j\rangle|^2=\sum_j\|T^*e_j\|^2=\|T^*\|_{HS}^2.$$
  Niech teraz $\{f_i\;:\;i\in I\}$ będzie konkurencyjna bazą ortonormalną $\Hh$. Wówczas
  $$\sum_i\|Te_i\|^2=\sum_i\sum_j|\langle Te_i, f_j\rangle|^2=\sum_i\sum_j|\langle e_i, T^*f_j\rangle|^2=\sum_j\sum_i|\langle e_i, T^*f_j\rangle|^2=\sum_j\|T^*f_j\|^2.$$
\end{proof}

\begin{remark}{}{}
  Niech $B_{HS}(\Hh)$ będzie zbiorem wszystkich operatorów Hilberta-Schmidta na $\Hh$. Wówczas $B_{HS}(\Hh)$ jest obustronnym $*$-ideałem w $B(\Hh)$.
\end{remark}

Zauważmy, że
$$\|T\|_{HS}^2=\sum_{i\in I}\langle Te_i, Te_i\rangle=\tr(T*T),$$
odważnym posunięciem jest więc zdefiniowanie iloczynu Hilberta-Schmidta jako
$$\langle T,S\rangle_{HS}:=\tr(S^*T)=\sum_{i\in I}\langle Te_i, Se_i\rangle.$$

\begin{remark}{}{}
  Przestrzeń operatorów Hilberta-Schmidta z iloczynem $\langle\cdot,\cdot\rangle_{HS}$ tworzy przestrzeń Hilberta.
\end{remark}

\begin{theorem}{}{}
  $T$ jest operatorem Hilberta-Schmidta wtedy i tylko wtedy gdy $|T|=\sqrt{T^*T}$ jest operatorem Hilberta-Schmidta.
\end{theorem}

\begin{definition}{norma Schattena}{}
  Niech $p\in[1,\infty)$ oraz $T\in B(\Hh)$ będzie operatorem ograniczonym. Wówczas definiujemy \buff{$p$-normę Schattena} jako
  $$\|T\|_p:=[\tr(|T|^p)]^{\frac{1}{p}}.$$
\end{definition}

Dla $p=2$ norma Schattena jest normą Hilberta-Schmidta, dla $p=1$ dostajemy normę śladową (nuklearną), a dla $p=\infty$ dostajemy $\|T\|_{\infty}=\sup_{\|\xi\|=1}\|T\xi\|$ zwykłą normę operatorową.

\begin{lemma}{}{}
  $$\|T\|_{HS}\geq \|T\|$$
\end{lemma}

\begin{proof}
  Dla dowolnego $\epsilon>0$ niech $\xi_0$, $\|\xi_0\|=1$, będzie wektorem takim, że 
  $$\|T\xi_0\|\geq \|T\|-\epsilon.$$
  Możemy takie $\xi_0$ rozszerzyć do bazy ortonormalnej $\{\xi_i\;:\;i\in I\}$ by dostać
  $$\|T\|_{HS}^2=\sum_i\|T\xi_i\|^2\geq\|T\xi_0\|^2\geq (\|T\|-\epsilon)^2.$$
  Przechodząc z $\epsilon\to 0$ dostajemy pożądaną nierówność.
\end{proof}

\begin{example}[m]
  \item Rozważmy $\Hh=L^2(X,\mu)$ z $\mu(X)<\infty$. Dla $k\in L^2(X\times X)$ rozważmy operator
    $$(T_kf)(x)=\int_Xk(x,y)f(y)d\mu(y).$$
    Jego norma Hilberta-Schmidta zgadza się z $L^2$ normą $k$ (ćwiczenie):
    $$\|T_k\|^2_{HS}=\int_X\int_X |k(x,y)|^2 d\mu(y)d\mu(x).$$
    \bigskip 

    Co więcej, jeśli $k_1, k_2\in L^2(X\times X)$, to zachodzi
    $$T_{k_1}T_{k_2}=T_{k_3},$$
    gdzie
    $$k_3(x,y)=\int_X k_1(x,z)k_2(z,y)d\mu(z).$$
  \item Dla $\phi\in C(G)$ określamy operator splotu 
    $$(T_\phi f)(x)=(f\ast\phi)(x)=\int_G f(y)\phi(y^{-1}x)dy$$
\end{example}

% ==============================================
%
% Twierdzenie von Neumanna-Bożejki-iiiiiiii
%
% \begin{theorem}{GNS konstrukcja}{}
%   se napisze
% \end{theorem}
%
%
% \begin{definition}{dodatnio określony operator}{}
%   Reprezentacja grupy $G$ operatorami liniowymi $\phi:G\to B(\Hh)$ nazywa sie dodatnio określoną, jeśli
%   $$(\forall\;x_i\in G,\xi_j\in\Hh)\;\sum_{i,j\leq N}\langle \phi(x_ix_j^{-1})\xi_i\;|\;\xi_j\rangle\geq0$$
% \end{definition}
%
% twierdzenie o reprezentacji funkcji dodatnio określonej na grupie dyskretnej: jeśli $\pi:G\to U(\Hh)$ jest unitarna reprezentacja, to $\phi(x)=\langle \pi(x)\xi_0\;|\;\xi_0\rangle$ dla $\|\xi_0\|=1$ jest jedyna
%
% GNS to twierdzenie odwrotne
%
% Przykłady: $G=(\Z, +)$, $\mu$ - miara probabilistyczna na grp dualnej (torus)
%
% $\phi(n)=\frac{1}{2\pi}\int_0^{2\pi}e^{int}d\mu(t)$ - opis całkowity (Bochner)
%
% Przykład: Najmark: niech $T$ taki, że $\|T\|\leq1$, $T\in M_n(\C)$, $N=\infty, 1, 2,...$. Niech $\phi_T:\Z\to B(\Hh)$
% $$\phi_T(n)=\begin{cases}
%   T^n & n\geq 0\\ 
%   (T^*)^{-n}&n<0
% \end{cases}$$
%
% dla $T=z$ dla $z\in\C$, $|z|\leq 1$ to $\phi_T$ jest jądrem Poissant
%
%
%
